\section{Conclusion}
Neural Bellman Operators separate proposal, evaluation, improvement, and certification so that a computational policy can be used in an economic counterfactual. The substantive question is not whether a method contains a neural network, but whether its representation and its error account support the change in primitives being studied.

For rewards with fixed transitions, the policy-feature and optimistic-support geometry is established. We retain its complete finite-model implementation, action-sign exclusions, and repaired policies, while identifying precisely the part shared with scalar optimistic linear support. For changing transition laws, an endpoint chord alone is invalid. The cross-operator correction in Theorem \ref{thm:corrected_chord}, feasible Bernstein policy values, and statewise switching provide a continuum welfare guarantee that incorporates this missing effect. Approximate neural endpoints can enter through certified residual supersolutions rather than assumed exact evaluation.

In the stopped preference economy, operating benefits and liquidation levels reduce to a contract contrast. A risk-specific duration--effort ratio then characterizes regular movements of an indifference locus, and a matched restriction on deliberate preference adjustment identifies the option-value contribution to the ranking of risky positions. These results separate a mechanism from a sign reversal selected at two contracts. The numerical brackets, all-parameter welfare certificates, and first-action sign exclusions retain their different scopes.

The resource comparison establishes the importance of reusable structure on both sides of a computational comparison. A nonlinear convex critic improves on the fitted quadratic continuation in the recorded experiments, but an exact reusable parametric quadratic program is faster than learned deployment for the tested short quadratic economy. Neither that fact nor the mesh-only policy-bank ablation invalidates the transfer certificate. They determine what can be attributed to learning and what should instead be attributed to policy evaluation, structural optimization, and parameter reuse.

The continuous-time formulation, finite transition targets, neural proposals, recursive preferences, temporal selves, and persistent dynamic competition remain distinct parts of the framework. Keeping their assumptions and evidence explicit permits new representations and counterfactuals without changing the economic meaning of the reported welfare guarantee.
